\newpage
\section*{四、实验探讨与思考题}


\subsection*{1. 实验小结}

本次对五种单色光各测三次截止电压，得到 $h=(6.52\pm0.15)\times10^{-34}\,\mathrm{J\cdot s}$，相对公认值的误差为1.57\%，其中不确定度仅为拟合标准不确定度。截止电压随频率近似线性增大，与爱因斯坦光电效应方程相符。由截距估算的逸出功为 $(2.00\pm0.06)$ eV，红限频率为 $(4.91\pm0.05)\times10^{14}$ Hz。

伏安曲线显示电流随正向电压增加而逐渐趋缓；拓展数据较好支持电流与孔径面积的关系，而距离实验存在偏离严格反平方关系的迹象。20 V下尚未明显饱和，是实验时的不恰当处理。



\subsection*{2. 测定普朗克常数的关键是什么？怎样根据光电管的特性曲线选择适宜的测定遏止电压的方法？}

关键是准确测定各频率的遏止电压，由电压—频率拟合斜率求 $h=ek$。背景和反向电流较小时采用零电流法；背景明显时先测量并扣除背景；反向电流较大且迅速饱和时，可用拐点法。


\subsection*{3. 从遏止电压与入射光频率的关系曲线中，你能确定阴极材料的逸出功吗？}

理想情况下，拟合截距 $b=-W/e$，可得 $W=-eb$，本次为 $(2.00\pm0.06)$ eV。接触电势和零点偏移会影响截距，未校正时只能作为逸出功的估计。


\subsection*{4. 暗电流会怎样影响遏止电压的测量？实验中应怎样检查并减小这种影响？}

暗电流使总电流零点偏离光电流零点，造成遏止电压偏差。应充分预热、换档调零，在遮光条件下记录一个稳定的暗电流值，检查本底是否接近零，待示数稳定后测量截止电压。



